\section{Introduction}
\label{introduction}
In 2015, the member states of the United Nations (UN) adopted the Sustainable Development Goals (SDGs) \cite{UN_SDGs_2024}, consisting of 17 global goals that, when interconnected, form a plan to alleviate poverty and protect the planet, aiming for a prosperous future for its inhabitants by 2030. In this regard, and in the current context of the fight against climate change, the SDGs play an important role through several goals that directly address environmental issues, such as SDG 7 (clean energy), SDG 11 (sustainable cities and communities), SDG 12 (responsible consumption), SDG 13 (climate action), and SDGs 14-15 (ecosystem protection). Cities, which concentrate a large part of the population and consumption, have a key role in this fight by promoting responsible waste management and environmental awareness among their inhabitants \cite{mi2019cities}. Of all the waste generated by the population, used cooking oil (UCO) plays a dual role: if its treatment is neglected, it can cause severe environmental damage; if recycled, it becomes an ally of sustainability, helping meet the SDGs. 

Globally, annual UCO production is estimated to be 41 to 67 million tons \cite{kumar2025valorization}. In Europe, production is estimated to be between 1 and 2.5 million tons annually \cite{mannu2025sustainable}. Most of this production comes from households and the hospitality sector, with cities contributing the most due to their high population density \cite{wan2019waste}. Most UCO is disposed of down the sink, potentially harming aquatic life by creating a surface layer that prevents oxygen exchange in the water \cite{tamothran2022degradation}, or it is mixed with solid waste sent to landfills. Moreover, some estimates claim that just 1 liter of oil can contaminate up to 1 million liters of water \cite{kumar2015transformation}. In addition, other related problems, such as foul odors and increased water treatment costs, result from poor UCO disposal \cite{orjuela2020green}.

However, UCO can become a valuable raw material for producing second-generation biodiesel when disposed of correctly (e.g., by taking it to a collection point). As is known, biodiesel is a biodegradable fuel that emits fewer polluting gases than fossil fuels, so promoting UCO circular economy can be a key factor in reversing climate change \cite{firoz2017review}. In this regard, the European Commission has also recognized UCO's dual potential, identifying it as a priority waste stream that could reduce dependence on fossil fuels  \cite{tamothran2022degradation}.

In this context, recycling used cooking oil is a concrete action that contributes to the SDGs: it prevents water and soil pollution, encourages responsible and circular consumption habits, and enables the generation of clean energy, such as biodiesel. As a result, many cities worldwide seeking to meet climate and sustainability goals are increasingly integrating the principles of the circular economy into their urban planning programs, exemplifying how circular approaches to waste treatment can address current environmental problems \cite{petit2018circular}. Given their governance capacity, the effective treatment of UCO is an imperative measure, recognizing it not as a simple disposal problem but as a resource stream that must be reintegrated. 

However, designing effective recycling systems requires consideration of several factors. Among these, and perhaps the key factor, is intelligent planning of the most suitable location for UCO collection points. The need to strategically define collection points is in line with classic operational research problems related to the location of infrastructure, providing practical frameworks for optimizing the location of facilities 
based on various criteria \cite{gallego2024metaheuristic,hale2003location,karpova2025multi,vecina2025enhancing}. For example, positioning strategies can be adopted that consider citizens' convenience in disposing of used oils and fats, increasing accessibility and, consequently, the recycling rate. 

Traditionally, location problems have been modeled using classic formulations such as the Facility Location Problem (FLP) and the Maximal Covering Location Problem (MCLP) derivatives. On the one hand, FLP generally focuses on minimizing costs, whether those associated with facility establishment or operation \cite{adeleke2020facility}. On the other hand, MCLP focuses on maximizing demand coverage while restricting it to a limited number of facilities \cite{murray2016maximal}. Although both formulations can be expanded to account for additional criteria, their use has focused on a single objective, often resulting in somewhat simplistic approaches to addressing complex balance problems in public infrastructure planning, where cost, coverage, and equity criteria converge. This commitment, when deciding on the location of facilities, is not new and was already recognized as a fundamental challenge in urban infrastructure planning in the work of Morrill et al. \cite{morrill1977efficiency}, who demonstrated that approaches focused on minimizing costs often result in unequal distributions of services. In FLPs, inequality means equitable distribution of infrastructure across a geographical area, ensuring that service users have reasonable access regardless of their location or socioeconomic status \cite{morrill1977efficiency,kunzmann1998planning}.

Building on the classical FLP and its coverage-oriented derivatives, such as the MCLP, we extend the model to incorporate multiple criteria relevant to urban sustainability. In particular, we propose an extension of the FLP to simultaneously consider four key objectives: the number of containers (a proxy for installation cost), the population left without capacity, the spatial equity of that shortage across neighbourhoods, and the distance citizens have to walk. First, minimizing the installation cost allows local authorities to implement economically viable solutions. Second, optimizing population coverage ensures that as many citizens as possible have access to nearby collection points, thereby increasing the likelihood of participation and reducing improper disposal of UCO. Finally, incorporating spatial equity criteria ensures all citizens benefit from the installed infrastructure. This approximation allows for the design of a collection system that balances efficiency with social fairness. 

The proposed model has been implemented using a multi-objective evolutionary algorithm, enabling the simultaneous optimization of conflicting objectives, and evaluated in an extensive study that compares the survival strategies of NSGA-II and NSGA-III, formulations with and without the distance objective, and several crossover operators, and bounds the accessibility of the solutions with a Lagrangian relaxation of the $p$-median problem. We adapt it to the structure of the problem with a set representation, a discrete set of network sizes and a load-guided mutation, and show through an ablation study that these components drive its performance, while recombination, including two spatially structured crossovers, brings no improvement. We also show why Moran's index of spatial autocorrelation, a natural candidate for measuring spatial equity, is ill-posed as an optimization objective, and propose a measure on a fixed spatial support instead. The method is available as an open-source service integrated in a data space \footnote{\url{https://github.com/CitComAI-Hub/container-location-optimizer}}, which reads the container inventory of a city through NGSI-LD and the population from GHS-POP, so it can be applied to any city. The proposed algorithm is then tested across five Spanish cities using real-world data: Donostia, Valencia, Sevilla, Málaga, and Valladolid. The rest of the paper is organized as follows: Section \ref{literature_review} reviews the literature on facility location problems and waste management. Section \ref{problem_definition} defines the problem. Section \ref{ga_algorithm} describes the proposed algorithm. Section \ref{computational_experiments} presents the experimental results using real data from Spanish cities. Finally, Section \ref{conclusions} offers conclusions and describes future work.


 